% Math 215 Exam 3 notecard (Lessons 19-27). 6in x 5in sheet = two 3x5 faces (cut/fold on the centre line).
% Zero margins, lines set solid, hairline column rules. Build in place:  tectonic -X compile Unit3_Notecard_6x5.tex
\documentclass{article}
% ---------- the knobs ----------
\newcommand\bodysize{3.3}      % body type size (pt)
\newcommand\nlines{115}        % text lines per column: sets the line pitch (pitch = (360pt - 2 edgemargin - 1em)/(nlines-1))
\newcommand\sidemargin{0.3pt}  % left/right inset so edge glyphs are not clipped; top and bottom are 0
\newcommand\edgemargin{0pt}    % safety inset on all four sides if a cut ever clips text (1pt costs about 5 lines: re-fit after)
\newcommand\colgap{1.4pt}      % gap between columns, with a hairline rule down its middle
\newcommand\stamplines{9}      % height of the STAMP HERE block, in text lines
\newdimen\insetLR \insetLR=\dimexpr\sidemargin+\edgemargin\relax
\usepackage[paperwidth=6in,paperheight=5in,left=\insetLR,right=\insetLR,top=\edgemargin,bottom=\edgemargin,nohead,nofoot]{geometry}
\usepackage{fontspec}
\usepackage{unicode-math}
\usepackage{multicol}
\usepackage{xcolor}
\usepackage{eso-pic}
\usepackage{graphicx}

% ---------- fonts ----------
\newcommand\fontdir{/Users/joelr/Library/Fonts/}  % where the commercial fonts live
\setmainfont{Bell-Centennial-Std-Address_6527.ttf}[Path=\fontdir, AutoFakeBold=1.8, AutoFakeSlant=0.18]
\setmathfont{FiraMath-Regular.otf}[Scale=0.92]
\newfontfamily\symf{DejaVuSans.ttf}[Scale=0.85]
\newfontfamily\stampf{Helvetica 95 Black.ttf}[Path=\fontdir, FakeBold=2.5]  % FakeBold = outline stroke weight

% ---------- sizes ----------
% First baseline sits 0.78em below the top edge, last baseline 0.22em above the bottom edge, nlines baselines per column.
\makeatletter
\newdimen\em@body \em@body=\bodysize pt
\newdimen\pitch \pitch=\dimexpr(360pt-2\dimexpr\edgemargin\relax-\em@body)/(\nlines-1)\relax
\edef\pitchpt{\strip@pt\pitch}
\edef\scriptsz{\strip@pt\dimexpr0.745\em@body\relax}   % first-level scripts (integral limits live here)
\edef\sscriptsz{\strip@pt\dimexpr0.66\em@body\relax}
\DeclareMathSizes{\bodysize}{\bodysize}{\scriptsz}{\sscriptsz}
\makeatother

% ---------- page furniture: hairline column rules; the centre one (fold / cut line) a little heavier ----------
% Solid black, not gray and not dashed: a gray hairline halftones into dots, and a dash beside a line end reads as | or 1.
\makeatletter
\AddToShipoutPictureBG{\AtPageLowerLeft{\setlength\unitlength{1pt}%
  \linethickness{0.2pt}\color{black}%
  \@tempdima=\dimexpr(432pt-2\insetLR-3\dimexpr\colgap\relax)/4\relax   % column width
  \@tempdimb=\dimexpr\insetLR+\@tempdima+0.5\dimexpr\colgap\relax\relax    % centre of the first gap
  \put(\strip@pt\@tempdimb,0){\line(0,1){360}}%
  \put(\strip@pt\dimexpr432pt-\@tempdimb\relax,0){\line(0,1){360}}%
  \linethickness{0.35pt}\put(216,0){\line(0,1){360}}}}
\makeatother

% ---------- spacing: every line exactly one pitch, nothing stretches ----------
\setlength\columnsep{\colgap}
\setlength\parindent{0pt}
\setlength\parskip{0pt}
\pagestyle{empty}
\tolerance=3000 \emergencystretch=0.8em \hbadness=10000 \hfuzz=0.4pt \vbadness=10000 \vfuzz=2pt
\hyphenpenalty=30 \exhyphenpenalty=30 \binoppenalty=200 \relpenalty=100
\AtBeginDocument{%
  \thinmuskip=2mu \medmuskip=2mu plus 1mu minus 1mu \thickmuskip=3mu plus 2mu minus 1mu
  \lineskiplimit=-\maxdimen \lineskip=0pt \nulldelimiterspace=0pt \scriptspace=0.1pt
  \topskip=0.78\dimexpr\bodysize pt\relax \splittopskip=\topskip \maxdepth=2pt \splitmaxdepth=\maxdepth}

% ---------- scripts that stay inside their line slot ----------
% TeX's own positions put a subscript 0.33em below and a superscript 0.85em above the baseline, which collides once
% lines are set solid. These wrappers keep the native attachment but shift the (smashed) script back toward the line.
% In the body ^ and _ are active and mean \supx and \subx. Never stack them on one base: use \I or \ev for that.
\newcommand\supx[1]{\sp{\smash{\lower0.085em\hbox{$\scriptstyle#1$}}}}   % single superscript
\newcommand\subx[1]{\sb{\smash{\raise0.20em\hbox{$\scriptstyle#1$}}}}    % single subscript
\newcommand\supy[1]{\sp{\smash{\lower0.02em\hbox{$\scriptstyle#1$}}}}    % stacked pair: upper
\newcommand\suby[1]{\sb{\smash{\raise0.10em\hbox{$\scriptstyle#1$}}}}    % stacked pair: lower
\newcommand\cx{\mathord{\symup{\char"5E}}}                               % caret for exponents that would not fit: e^(2x)
\newcommand\I[2]{{\smash{\int}}\suby{#1}\supy{#2}}                        % integral with stacked limits (no powers inside)
\newcommand\Ia[2]{{\smash{\int}}[#1{\to}#2]}                              % integral with limits written [a -> b] (limits with powers)
\newcommand\J{{\smash{\int}}}                                             % integral sign alone
\newcommand\II[1]{{\smash{\iint}}\subx{#1}}                               % double integral over a region
\newcommand\III[1]{{\smash{\iiint}}\subx{#1}}                             % triple integral over a solid
\newcommand\ev[2]{\suby{#1}\supy{#2}}                                     % evaluation limits: [F]\ev{a}{b}
\newcommand\pr{{\smash{\raise0.30em\hbox{$\scriptstyle\prime$}}}}         % prime (math ' is mapped to this in the body)

% ---------- markup ----------
\newcommand\vb[1]{\symbf{#1}}                         % bold vector
\newcommand\nm[1]{\mathopen{\smash{\lVert}}#1\mathclose{\smash{\rVert}}}     % norm
\newcommand\ip[1]{\mathopen{\smash{\langle}}#1\mathclose{\smash{\rangle}}}   % <a,b,c>
\newcommand\abs[1]{\mathopen{\smash{\lvert}}#1\mathclose{\smash{\rvert}}}    % |x|
\newcommand\cm{,\allowbreak\,}                        % comma that lets a long vector break across lines
\newcommand\sh[1]{{\setlength\fboxsep{0.18pt}\colorbox{black!22}{\textbf{#1}}}}  % equation-sheet tag
\newcommand\hd[1]{\par\noindent{\setlength\fboxsep{0pt}\smash{\colorbox{black}{\color{white}\bfseries
  \vrule height0.76em depth0.20em width0pt\,#1\,}}}\ }                           % run-in section label, one line slot tall
\newcommand\pt[1]{\textbf{{\symf ▸}#1}}       % problem type (run-in); follow it with a space or a colon
\newcommand\ex{\textit{Ex}\ }                  % worked example
\newcommand\how{\textbf{How:}\ }
\newcommand\warn{{\symf ⚠}}                    % pitfall
\newcommand\shc{{\symf ★}}                      % shortcut / extra
\newcommand\ok{{\symf ✓}}                       % checked
\newcommand\rt{\smash{\surd}}                  % prefix radical: \rt2, \rt(x^2+y^2)
\newcommand\then{\ensuremath{\Rightarrow}}
\newcommand\pd{\partial}
\newcommand\gr{\nabla}
\newcommand\av{\symup{avg}}

\begin{document}
\fontsize{\bodysize}{\pitchpt}\selectfont
\catcode`\^=\active \catcode`\_=\active \let^\supx \let_\subx
\begingroup\lccode`\~=`\' \lowercase{\endgroup\let~}\pr
\begin{multicols*}{4}

\hd{MATH 215 EXAM 3 · L19–L27 · HOW TO READ THE TAGS}
A gray tag points at the allowed equation sheet and stands in for the formula: read the formula there, read here what it does, what every letter means and how to use it. \sh{7M5} = \textbf{page 7, M}iddle (full-width) rows in the lower half, \textbf{5}th one down. \sh{7L4}, \sh{7R4} = p7 \textbf{L}eft / \textbf{R}ight column, 4th formula from the top. \sh{6UL2}, \sh{6UR2} = p6, the block under the heading ``\textbf{U}nit 3 Equations'', left / right column, 2nd formula. \sh{5C1}–\sh{5C6} = p5, the six centered coordinate-conversion rows. \sh{4L7} = p4 left column, 7th row. \sh{3L7}, \sh{3R7} = p3 table, row 7 under the header: left cell (derivative), right cell (integral). \sh{p1} = trig identities (by name), \sh{p2} = unit circle. \shc = shortcut, \warn = pitfall, \ex = worked example, \then = ``so / gives''. $f_x=\pd f/\pd x$ = partial derivative (see TOOLKIT, last section). Exponents that would not fit on the line use a caret: $e\cx(2x)$ = $e$ to the power $2x$. $\Ia{a}{b}$ = integral from $a$ to $b$ (same as $\I{a}{b}$).

\hd{L19 CHAIN RULE \& IMPLICIT DERIVATIVES}
Idea: when the inputs of $f$ are themselves functions, the rate of change of $f$ is the sum, over every input, of (how sensitive $f$ is to that input) × (how fast that input moves).
\pt{One parameter} \sh{6UL1}: $z=f(x,y)$ with $x=g(t)$, $y=h(t)$, so $z$ is really a function of $t$ alone and the answer is an ordinary derivative $dz/dt$. $\pd z/\pd x$, $\pd z/\pd y$ = partials of the outer function (the other input frozen); $dx/dt$, $dy/dt$ = ordinary derivatives of the inner functions. \how 1) $f_x$, $f_y$; 2) $x'(t)$, $y'(t)$; 3) multiply in pairs, add; 4) replace $x$, $y$ by their $t$-formulas (``in terms of $t$''). \ex $f=4x^2+3y^2$, $x=\sin t$, $y=\cos t$: $(8x)(\cos t)+(6y)(-\sin t)=8\sin t\cos t-6\sin t\cos t=2\sin t\cos t$. \ex $f=\rt(x^2-y^2)$, $x=e\cx(2t)$, $y=e\cx(-t)$: $f_x=x/\rt(x^2-y^2)$, $f_y=-y/\rt(x^2-y^2)$, $x'=2e\cx(2t)$, $y'=-e\cx(-t)$ \then\ $(2e\cx(4t)+e\cx(-2t))/\rt(e\cx(4t)-e\cx(-2t))$. \ex $f=e\cx(8xy+5)$, $x=t^{1/3}$, $y=t^2$: $f_x=8ye\cx(..)$, $f_y=8xe\cx(..)$, $x'=(1/3)t^{-2/3}$, $y'=2t$ \then\ $(8/3)t^{4/3}e\cx(..)+16t^{4/3}e\cx(..)=(56/3)t^{4/3}e\cx(8t^{7/3}+5)$. \ex (Quiz 5) $f=e\cx(9xy+8)$, $x=t^{1/4}$, $y=t^2$ \then\ $(9/4+18)t^{5/4}e\cx(..)=(81/4)t^{5/4}e\cx(9t^{9/4}+8)$. \shc Check: substitute first, $f=e\cx(8t^{7/3}+5)$, then the one-variable chain rule gives the same thing. \shc Tree: $z$ branches to $x$ and $y$, each branches to $t$; multiply along each path, add the paths.
\pt{Two parameters} \sh{6UL2}\,\sh{6UL3}: $x=g(u,v)$, $y=h(u,v)$, so $z$ depends on $u$ and $v$ and the answers are partials: $\pd z/\pd u$ (hold $v$ fixed), $\pd z/\pd v$ (hold $u$ fixed). Same recipe, using the inner partials for the letter asked. \ex $w=6x^4-8y^4$, $x=s-8t$, $y=s+7t$: $w_t=(24x^3)(-8)+(-32y^3)(7)=-192(s-8t)^3-224(s+7t)^3$; $w_s=24(s-8t)^3-32(s+7t)^3$. \ex (Quiz 5) $w=9x^2+8y$, $x=-3s+4t$, $y=4s+5t$: $w_t=(18x)(4)+(8)(5)=72x+40=-216s+288t+40$.
\pt{More variables} (generalized rule, not on the sheet): $\pd w/\pd t_j=\sum_i(\pd w/\pd x_i)(\pd x_i/\pd t_j)$ = one term per middle variable $x_i$; a middle variable that does not contain $t_j$ adds 0. \ex $w=-2x^2-2y^2+4z^2$, $x=3\cos t$, $y=2\sin t$, $z=5e^t$: $(-4x)(-3\sin t)+(-4y)(2\cos t)+(8z)(5e^t)=36\sin t\cos t-16\sin t\cos t+200e\cx(2t)=20\sin t\cos t+200e\cx(2t)$. \ex $f=2z(4x+8y)$, $x=-2v+7w$, $y=w$, $z=8uv$: $f_w=(8z)(7)+(16z)(1)+(8x+16y)(0)=72z=576uv$.
\pt{Implicit differentiation} (the equation is not solved for $y$ or $z$, but you want its slope anyway). \textbf{Formula way}: move everything to one side so it reads $F=0$; then $dy/dx=-F_x/F_y$ \sh{6UL4}, and with three variables $\pd z/\pd x=-F_x/F_z$ \sh{6UR1}, $\pd z/\pd y=-F_y/F_z$ \sh{6UR2} (the sheet names $F$ ``$f$''). Why: $F$ stays 0 along the curve, so the chain rule gives $F_x+F_y\,y'=0$. $F_x$, $F_y$, $F_z$ = ordinary partials of that one-sided expression with every other letter frozen (here $y$ and $z$ are just letters, not functions). \warn\ (1) minus sign out front; (2) it is upside down from how it reads: the letter you differentiate \textit{with respect to} goes on TOP, the dependent letter ($y$ or $z$) on the BOTTOM; (3) the bottom must not be 0; (4) a constant on the right ($\dots=-4$) disappears, so you may differentiate the left side as it stands. \ex $x^2+3y^2+4y-4=0$: $F_x=2x$, $F_y=6y+4$ \then\ $dy/dx=-2x/(6y+4)=-x/(3y+2)$. \ex $-5x+5y^4+8y^2-10=0$ \then\ $dy/dx=-(-5)/(20y^3+16y)=5/(20y^3+16y)$. \ex (Quiz 5) $4x^4z^3+x^2y^3+y^4+3yz^3-z^3=1$: $F_x=16x^3z^3+2xy^3$, $F_z=12x^4z^2+9yz^2-3z^2$ \then\ $\pd z/\pd x=-(16x^3z^3+2xy^3)/(12x^4z^2+9yz^2-3z^2)$. \ex $-5x^4z^4+2x^4-y^4z^3-5y^2-2z^2=-4$: $F_z=-20x^4z^3-3y^4z^2-4z$ \then\ $z_x=-(-20x^3z^4+8x^3)/F_z$, $z_y=-(-4y^3z^3-10y)/F_z$. \textbf{Calc 1 way} (same answer; use it when told to differentiate both sides): take $d/dx$ of both sides treating $y$ as a hidden function $y(x)$, so every $y$-piece picks up a factor $y'$ from the chain rule: $(y^3)'=3y^2y'$, $(\sin y)'=y'\cos y$; a mixed term needs the product rule: $(x^2y^3)'=2xy^3+x^2\cdot3y^2y'$; then move every $y'$ term to one side, factor $y'$ out, divide. \ex $x^2+3y^2+4y-4=0$: $2x+6yy'+4y'=0$ \then\ $y'(6y+4)=-2x$ \then\ $y'=-x/(3y+2)$ \ok. \ex $x^3+y^3=6xy$: $3x^2+3y^2y'=6y+6xy'$ \then\ $y'(3y^2-6x)=6y-3x^2$ \then\ $y'=(2y-x^2)/(y^2-2x)$; at $(3,3)$: $-3/3=-1$. \ex Slope at a point by the formula: $x^2y+y^3=10$ at $(1,2)$: $dy/dx=-2xy/(x^2+3y^2)=-4/13$. For $z(x,y)$: differentiate in $x$ with $y$ frozen and give every $z$-piece a $z_x$: $x^2z+z^3=y$ \then\ $2xz+x^2z_x+3z^2z_x=0$ \then\ $z_x=-2xz/(x^2+3z^2)$. Tangent line at $(a,b)$: $y-b=y'(a,b)\,(x-a)$ (plug the point in last). Horizontal tangent: points of the curve with $F_x=0$, $F_y\ne0$.

\hd{L20 DIRECTIONAL DERIVATIVE \& GRADIENT}
\pt{Gradient} $\gr f$ \sh{6UR4} (2 variables), \sh{7L1} (3 variables): the vector of first partials, $\ip{f_x,f_y}$ or $\ip{f_x,f_y,f_z}$ ($\vb i,\vb j,\vb k$ = unit vectors along $x,y,z$; also written grad $f$, ``del $f$''). Meaning: it points in the direction in which $f$ rises fastest (straight uphill) and its length is that steepest slope. ``At a point'' = plug the point into every component. \ex $-7x^3y+2xy^2$ \then\ $\ip{-21x^2y+2y^2\cm-7x^3+4xy}$. \ex $-7x\ln(2y)$ \then\ $\ip{-7\ln(2y)\cm-7x/y}$. \ex $\sin9x\sin9y$ at $(\pi/4,\pi/3)$: $\ip{9\cos9x\sin9y\cm9\sin9x\cos9y}$ with $9\pi/4\to\pi/4$, $\sin3\pi=0$, $\cos3\pi=-1$ \then\ $\ip{0\cm-9\rt2/2}$. \ex $5x^2-2xy+y^2-4yz+z^2+3xz$ \then\ $\ip{10x-2y+3z\cm-2x+2y-4z\cm3x-4y+2z}$. \ex (Quiz 5) $-7x^3yz^2+8xy^2z^2-7xyz^2$ \then\ $\ip{-21x^2yz^2+8y^2z^2-7yz^2\cm-7x^3z^2+16xyz^2-7xz^2\cm-14x^3yz+16xy^2z-14xyz}$.
\pt{Directional derivative} $D_{\vb u}f$ = the slope of $f$ (change in $f$ per unit distance) as you leave the point heading in direction $\vb u$. $\vb u$ must be a UNIT vector (length 1). \textbf{Dot form} \sh{6UR5} (2D), \sh{7L2} (3D): gradient dotted with $\vb u$. \textbf{Angle form} \sh{6UR3}: $\theta$ = direction angle, counterclockwise from the $+x$-axis, i.e. $\vb u=\ip{\cos\theta,\sin\theta}$ (already unit). \how 1) $\gr f$; 2) plug in the point; 3) make the direction unit: $\vb u=\vb v/\nm{\vb v}$, $\nm{\ip{a,b}}=\rt(a^2+b^2)$ \sh{4L11}\,\sh{4L12}; ``from $P$ toward $Q$'' means $\vb v=Q-P$; 4) dot \sh{4L16}: multiply matching components, add \then\ one number. \warn\ normalize the direction, not the gradient; evaluate the gradient before dotting. \shc Compute $(\gr f\cdot\vb v)/\nm{\vb v}$: no fractions until the end; then rationalize $a/\rt b=a\rt b/b$. \ex (Quiz 5) $2x^3y^2-3x^2y-7xy^3$, $P(-2,1)$, $\vb v=\ip{4,-2}$: $\gr f=\ip{6x^2y^2-6xy-7y^3\cm4x^3y-3x^2-21xy^2}=\ip{29,-2}$, $\nm{\vb v}=2\rt5$ \then\ $(116+4)/(2\rt5)=60/\rt5=12\rt5$. \ex (Quiz 5) $8x^2\sin4y$, $P(-2,\pi/2)$, $\vb v=\ip{1,2}$: $\gr f=\ip{16x\sin4y\cm32x^2\cos4y}=\ip{0,128}$ \then\ $256/\rt5=256\rt5/5$. \ex $x^3y^3-2x^2y$, $P(-2,-4)$, $\vb v=\ip{-3,2}$: $\gr f=\ip{3x^2y^3-4xy\cm3x^3y^2-2x^2}=\ip{-800,-392}$ \then\ $(2400-784)/\rt13=1616\rt13/13$. \ex $-\sin x\cos4y$ from $P(0,0)$ toward $Q(2\pi/3,\pi)$: $\vb v=\ip{2\pi/3,\pi}\parallel\ip{2,3}$, $\gr f(0,0)=\ip{-1,0}$ \then\ $-2/\rt13=-2\rt13/13$. \ex 3D: $5x^2-2xy+y^2-4yz+z^2+3xz$ (gradient found above) at $(-1,-2,3)$, $\vb v=\ip{-1,2,2}$: $\gr f=\ip{3,-14,11}$, $\nm{\vb v}=3$ \then\ $(-3-28+22)/3=-3$. \ex angle form: $3x^2y^3+3x^2y+8xy-1$, $\theta=0$ \then\ $f_x\cdot1+f_y\cdot0=6xy^3+6xy+8y$. \ex $9y^4e\cx(2x+4y)$ at $(2,5)$, $\theta=\pi$ \then\ $-f_x=-18y^4e\cx(2x+4y)=-11250e^{24}$. \shc $\theta=0$: $f_x$; $\pi/2$: $f_y$; $\pi$: $-f_x$; $3\pi/2$: $-f_y$. Limit definition (not on the sheet): $D_{\vb u}f(a,b)=\lim_{h\to0}[f(a+h\cos\theta\cm b+h\sin\theta)-f(a,b)]/h$. A 3D unit vector is $\ip{\cos\alpha,\cos\beta,\cos\gamma}$ (direction cosines = cosines of its angles with the $x,y,z$ axes; their squares add to 1).
\pt{Steepest directions.} $D_{\vb u}f=\nm{\gr f}\cos\phi$, $\phi$ = angle between $\gr f$ and $\vb u$. So the \textbf{largest} rate is $\nm{\gr f}$, in the direction $\gr f/\nm{\gr f}$; the \textbf{smallest} (most negative) is $-\nm{\gr f}$, in the direction $-\gr f/\nm{\gr f}$; the rate is 0 perpendicular to $\gr f$ (along the level curve); if $\gr f=\vb0$ every directional derivative is 0. \ex (Quiz 5) direction of minimum for $-2x^3y^3-x^3y^2+4xy^2$ at $(-2,-1)$: $\gr f=\ip{-6x^2y^3-3x^2y^2+4y^2\cm-6x^3y^2-2x^3y+8xy}=\ip{16,48}$ \then\ $\vb u=\ip{-\rt10/10\cm-3\rt10/10}$. \shc Pull the common factor out before normalizing: $\ip{16,48}=16\ip{1,3}$ \then\ unit vector $\ip{1,3}/\rt10$. \ex direction of maximum for $3x^3y^3-3x^2y^3-xy$ at $(-1,2)$: $\gr f=\ip{9x^2y^3-6xy^3-y\cm9x^3y^2-9x^2y^2-x}=\ip{118,-71}$ \then\ $\vb u=\ip{118,-71}/\rt18965$, and the maximum rate is $\rt18965$.
\pt{Level curves: normal and tangent.} $\gr f(x_0,y_0)$ is perpendicular (normal) to the level curve $f=k$ through that point. Unit normal: $\gr f/\nm{\gr f}$. Unit tangent = the normal turned 90°: swap the two components and negate one, $\pm\ip{-f_y,f_x}/\nm{\gr f}$ (either sign is a correct answer). \ex $x^2-2y$ at $(2,1)$: $\gr f=\ip{2x,-2}=\ip{4,-2}$ \then\ unit normal $\ip{2\rt5/5\cm-\rt5/5}$. \ex (Quiz 5) $5y^2/(5x-1)$ at $(2,1)$: $f_x=-25y^2/(5x-1)^2=-25/81$, $f_y=10y/(5x-1)=10/9$; \shc clear the fractions (×81): $\parallel\ip{-25,90}\parallel\ip{-5,18}$, length $\rt349$ \then\ tangent $\ip{18,5}/\rt349=\ip{18\rt349/349\cm5\rt349/349}$. \ex $-2x^2/(y+2)$ at $(-1,0)$: $\gr f=\ip{-4x/(y+2)\cm2x^2/(y+2)^2}=\ip{2,1/2}\parallel\ip{4,1}$ \then\ tangent $\ip{-1,4}/\rt17=\ip{-\rt17/17\cm4\rt17/17}$. \shc Surface $F(x,y,z)=k$: $\gr F$ is normal to it, so the tangent plane is $F_x(x-x_0)+F_y(y-y_0)+F_z(z-z_0)=0$ \sh{4R16}.

\hd{L21 MAXIMA \& MINIMA}
\pt{Critical points} \sh{7L3}: the points $(x_0,y_0)$ where both first partials are 0 (flat tangent plane), or where one of them does not exist. Local max at $(x_0,y_0)$: $f(x_0,y_0)\ge f(x,y)$ for every point of some small disk around it (local min: $\le$); absolute = over the whole region. Every local max or min happens at a critical point (Fermat), but a critical point need not be either. \how set $f_x=0$ and $f_y=0$ and solve them together (TOOLKIT: systems). \ex (Quiz 5) $x^2+2xy-8y^2-4x+4y-8$: $2x+2y-4=0$ and $2x-16y+4=0$; subtract \then\ $18y=8$, $y=4/9$, $x=2-y=14/9$ \then\ $(14/9,4/9)$. \ex $x^2+2xy-4y^2+4x-6y+4$: $2x+2y+4=0$, $2x-8y-6=0$ \then\ $(-1,-1)$. \ex $x^2+2xy-6y^2+2x-8y-2$ \then\ $(-2/7,-5/7)$. \ex $(4x^2+9y^2+8)^{4/3}$: $f_x=(4/3)(..)^{1/3}(8x)$, $f_y=(4/3)(..)^{1/3}(18y)$; the bracket is at least 8, never 0 \then\ only $(0,0)$. Nonlinear systems: factor, split into cases, and pair each case with the other equation.
\pt{Second derivative test} \sh{7L4}: at a critical point compute the discriminant $D$ = (the two pure second partials multiplied) minus (the mixed partial squared). $f_{xx}$ = how the surface bends in the $x$-direction; $f_{xy}=(f_x)_y$. \textbf{$D>0$ and $f_{xx}>0$: local min} (bowl); \textbf{$D>0$ and $f_{xx}<0$: local max} (dome); \textbf{$D<0$: saddle} (up one way, down another: no extremum); \textbf{$D=0$: the test says nothing}. \warn\ evaluate $D$ separately at each critical point; the ``local max/min value'' is $f$ at the point. \ex $x^3/3+y^2+2xy-6x-3y+4$: $f_x=x^2+2y-6$, $f_y=2y+2x-3$ \then\ $2y=3-2x$ \then\ $x^2-2x-3=0$, $x=3$ or $-1$ \then\ $(3,-3/2)$ and $(-1,5/2)$. $f_{xx}=2x$, $f_{yy}=2$, $f_{xy}=2$ \then\ $D=4x-4$. At $(3,-3/2)$: $D=8>0$, $f_{xx}=6>0$ \then\ local min, value $-29/4$. At $(-1,5/2)$: $D=-8<0$ \then\ saddle. \ex (Quiz 5) $-6x^2-3x+2y^2+6y$ at $(-1/4,-3/2)$: $D=(-12)(4)-0=-48$ \then\ saddle. \ex $9x^2-7x+y^2-2y$ at $(7/18,1)$: $D=(18)(2)-0=36$, $f_{xx}=18$ \then\ local min. \shc Quadratic $ax^2+bxy+cy^2+\dots$: $D=4ac-b^2$ and $f_{xx}=2a$ at every point (nothing to plug in). With no $xy$ term: $x^2$ and $y^2$ coefficients of opposite sign \then\ saddle; both positive \then\ min; both negative \then\ max. \shc Completing the square can show the answer at once: $(3x+4)^2+(3y-5)^2\ge0$ \then\ absolute min 0 at $(-4/3,5/3)$; $x^2+2xy+2x+2y^2+6y+8=(x+y+1)^2+(y+2)^2+3$ \then\ min 3 at $(1,-2)$. \ex $x^3+y^3-3xy$: $3x^2=3y$, $3y^2=3x$ \then\ $(0,0)$: $D=-9$, saddle; $(1,1)$: $D=27$, $f_{xx}=6$, local min, value $-1$.
\pt{Absolute max and min on a closed, bounded region} (closed = contains its boundary; bounded = fits inside some disk; Extreme Value Theorem: a continuous $f$ on such a region reaches both). \how 1) critical points inside the region; 2) boundary: on each edge substitute its equation ($x=c$, $y=c$ or $y=g(x)$) to get a one-variable function, set its derivative to 0, keep the points that lie on the edge; 3) the corners / endpoints; 4) evaluate $f$ at every candidate: largest = absolute max, smallest = absolute min (no $D$ test needed). \ex $x^2-2xy+4y^2-4x-2y+24$ on $0\le x\le4$, $0\le y\le2$. Inside: $2x-2y-4=0$, $-2x+8y-2=0$ \then\ $(3,1)$, $f=17$. Edges: $y=0$: $x^2-4x+24$ \then\ $x=2$, $f=20$; $y=2$: $x^2-8x+36$ \then\ $x=4$, $f=20$; $x=0$: $4y^2-2y+24$ \then\ $y=1/4$, $f=95/4$; $x=4$: $4y^2-10y+24$ \then\ $y=5/4$, $f=71/4$. Corners: $f(0,0)=24$, $f(4,0)=24$, $f(0,2)=36$, $f(4,2)=20$. \then\ absolute max 36 at $(0,2)$, absolute min 17 at $(3,1)$. Circular boundary: use Lagrange (L22) or $x=R\cos t$, $y=R\sin t$. No critical point inside \then\ both extremes are on the boundary.

\hd{L22 LAGRANGE MULTIPLIERS}
\pt{One constraint} \sh{7L5}: the max/min of $f$ (the objective: what you are optimizing) when you may only stand on the curve $g=0$ (the constraint, with everything moved to one side: $x+2y=7$ becomes $g=x+2y-7$; the constant never affects $\gr g$). At the best point a level curve of $f$ just touches the constraint curve, so the two gradients are parallel; $\lambda$ (the multiplier) is the unknown stretch factor between them (it may be negative, and you rarely need its value). $(x_0,y_0)$ = the point you are solving for. \how 1) name $f$ and $g$; 2) write $f_x=\lambda g_x$, $f_y=\lambda g_y$ (and $f_z=\lambda g_z$), plus the constraint itself; 3) get rid of $\lambda$: solve two equations for $\lambda$ and set them equal, or divide one by the other \then\ a link between $x$ and $y$; 4) put the link into the constraint; 5) plug every candidate into $f$: biggest = max, smallest = min. \warn\ the answer is a value of $f$ (or the point), never $\lambda$. \warn\ never divide by something that might be 0: factor instead, $y(6-2\lambda)=0$ \then\ $y=0$ OR $\lambda=3$, and follow both branches. \ex min of $x^2+4y^2-2x+8y$ on $x+2y=7$: $2x-2=\lambda$, $8y+8=2\lambda$ \then\ $8y+8=4x-4$, $y=(x-3)/2$; constraint: $2x-3=7$ \then\ $(5,1)$, $f=27$ (a bowl along a line has a lowest point only). \shc Line constraint: substitute instead ($x=7-2y$ \then\ $8y^2-16y+35$) and take the vertex $y=-B/(2A)=1$. \ex min of $7x+4y$ on $x^2+y^2=32$: $7=2\lambda x$, $4=2\lambda y$; divide \then\ $y=(4/7)x$ \then\ $(65/49)x^2=32$, $x=\pm28\rt130/65$; $f=(65/7)x$ \then\ min $-4\rt130$ (max $+4\rt130$). \shc Linear $f=ax+by$ on the circle $x^2+y^2=R^2$: max and min are $\pm R\rt(a^2+b^2)$, at the points $\pm R\ip{a,b}/\rt(a^2+b^2)$. (Quiz 5) $-9x-8y$ on $x^2+y^2=8$: min $=-\rt8\,\rt145=-2\rt290$. \ex Three variables: min of $x^2+y^2+z^2$ on $x+y+z=1$: $2x=2y=2z=\lambda$ \then\ $x=y=z=1/3$ \then\ $1/3$ (the squared distance from the origin to the plane; \shc for the plane $ax+by+cz=d$ it is $d^2/(a^2+b^2+c^2)$). \shc Distance problems: optimize the squared distance (no root), take the root at the end. \ex Zero branch: $2xy^2$ on $4x^2+y^2=9$: $2y^2=8\lambda x$, $4xy=2\lambda y$ \then\ $y(4x-2\lambda)=0$. Branch $y=0$: $x=\pm3/2$, $f=0$. Branch $\lambda=2x$: $y^2=8x^2$ \then\ $12x^2=9$, $x=\pm\rt3/2$, $y^2=6$, $f=12x=\pm6\rt3$ \then\ max $6\rt3$, min $-6\rt3$. \ex Word problem: open-top box from 24 ft$^2$ of cardboard, largest volume: $f=xyz$, $g=xy+2xz+2yz-24$: $yz=\lambda(y+2z)$, $xz=\lambda(x+2z)$, $xy=\lambda(2x+2y)$; multiply the first by $x$ and the second by $y$ (both left sides become $xyz$) \then\ $x=y$; second and third the same way \then\ $y=2z$; constraint: $12z^2=24$ \then\ $z=\rt2$, $x=y=2\rt2$, $V=8\rt2\approx11.31$ ft$^3$.
\pt{Two constraints} (not on the sheet): $\gr f=\lambda_1\gr g+\lambda_2\gr h$ together with $g=0$ and $h=0$ (three component equations + two constraints; unknowns $x,y,z,\lambda_1,\lambda_2$). \ex max of $x+2y+3z$ on $x-y+z=1$ and $x^2+y^2=1$: $1=\lambda_1+2\lambda_2x$, $2=-\lambda_1+2\lambda_2y$, $3=\lambda_1$ \then\ $x=-1/\lambda_2$, $y=5/(2\lambda_2)$; cylinder: $29/(4(\lambda_2)^2)=1$ \then\ $\lambda_2=\pm\rt29/2$; plane: $z=1-x+y$ \then\ $f=3\pm\rt29$ \then\ max $3+\rt29$.
\pt{Filled region} ($x^2+y^2\le R^2$ and the like): candidates = critical points in the interior ($\gr f=\vb0$, L21) plus Lagrange on the boundary; compare them all. \ex max of $-9x^2-4x+3y^2+5$ on $x^2+y^2\le25$. Inside: $-18x-4=0$, $6y=0$ \then\ $(-2/9,0)$, $f=49/9$. Boundary: $-18x-4=2\lambda x$, $6y=2\lambda y$ \then\ $y=0$: $(\pm5,0)$, $f=-240$ and $-200$; or $\lambda=3$: $x=-1/6$, $y^2=899/36$, $f=241/3$ \then\ max $241/3$ at $(-1/6,\pm\rt899/6)$ (min $-240$). \ex $-8x^2-5y^2+3y+1$ on $x^2+y^2\le49$: inside $(0,3/10)$, $f=29/20$ = max; boundary $x=0$: $f(0,\pm7)=-223,-265$; $\lambda=-8$: $y=-1/2$, $x^2=195/4$, $f=-1567/4$ = min. Lone candidate: test another constraint point to tell max from min.

\hd{L23 DOUBLE INTEGRALS OVER RECTANGLES}
\pt{Meaning.} $\II{R}f\,dA$ = the signed volume between the surface $z=f(x,y)$ and the region $R$ of the $xy$-plane (parts below the plane count negative) = the limit of the Riemann sum $\sum\sum f\,\Delta A$: chop $R$ into small rectangles of area $\Delta A=\Delta x\,\Delta y$ with $\Delta x=(b-a)/m$, $\Delta y=(d-c)/n$, and stand a box of height $f$(sample point) on each. $f$ is ``integrable'' when that limit exists; every continuous $f$ is (same for triple integrals). $dA$ = tiny piece of area = $dx\,dy$ = $dy\,dx$. $R=[a,b]\times[c,d]$ means $a\le x\le b$ and $c\le y\le d$ (\warn\ the first bracket is $x$).
\pt{Fubini, iterated integral} \sh{7M1}: on a rectangle the double integral is two ordinary integrals, in either order. $a,b$ = the $x$-limits; $c,d$ = the $y$-limits. \how work inside out. The inner differential names the variable that goes first and owns the inner limits; while integrating in $y$ every $x$ is a constant (a partial derivative run backwards); plug in the inner limits \then\ a function of the outer variable only; integrate that. \ex $\II{R}(2x^2-7x^2y+2y)\,dy\,dx$, $R=[0,2]\times[3,4]$: inner $[2x^2y-(7/2)x^2y^2+y^2]\ev{y=3}{4}=7-(45/2)x^2$; outer $[7x-(15/2)x^3]\ev{0}{2}=14-60=-46$ (negative: mostly below the plane). \ex $\II{R}3x^2\cos6y\,dy\,dx$, $R=[0,\pi/2]\times[\pi/4,1]$: inner $3x^2[(\sin6y)/6]\ev{\pi/4}{1}$ \sh{3R7} $=(x^2/2)(\sin6+1)$ since $\sin(3\pi/2)=-1$ ($\sin6$ = sine of 6 radians: leave it) \then\ $((1+\sin6)/2)(\pi/2)^3/3=\pi^3(1+\sin6)/48$.
\pt{Properties} (any region): the integral of a sum is the sum of the integrals; constants pull out; cut the region into pieces $S$ and $T$ and add the two integrals; $f\le g$ on $R$ \then\ $\II{R}f\le\II{R}g$; $m\le f\le M$ \then\ $m\cdot$Area $\le\II{R}f\le M\cdot$Area. \ex $\II{R}f=5$, $\II{R}g=5$ \then\ $\II{R}(-9f+5g)=-45+25=-20$. \shc \textbf{Product rule} (rectangle with number limits, integrand = ($x$-only)×($y$-only)): $\II{R}g(x)h(y)\,dA=(\I{a}{b}g\,dx)(\I{c}{d}h\,dy)$. \ex $\II{R}3x^2y^2\cos(6y^3)\,dA$, $-4\le x\le3$, $-\pi/6\le y\le\pi/6$: $(\I{-4}{3}3x^2\,dx)(\Ia{-\pi/6}{\pi/6}y^2\cos(6y^3)\,dy)=[x^3]\cdot[(1/18)\sin(6y^3)]$ ($u=6y^3$, see TOOLKIT) $=(91)(1/9)\sin(\pi^3/36)$. \ex $y^2e\cx(x-7)$ on $[0,1]\times[1,4]$: $(e\cx(-6)-e\cx(-7))(21)=21(e-1)/e^7$. \ex $y\rt(6x+7)$ on $[0,1]\times[0,4]$: $[(6x+7)^{3/2}/9]\ev{0}{1}\cdot8=(8/9)(13^{3/2}-7^{3/2})$. \shc A sum splits too: $\II{R}(g(x)+h(y))\,dA=(d-c)\I{a}{b}g\,dx+(b-a)\I{c}{d}h\,dy$.
\pt{Area} \sh{7R1}: integrate the function 1 (a slab of height 1 has volume = its base area). $[0,3]\times[0,2]$ \then\ $\I{0}{3}\I{0}{2}1\,dy\,dx=6$. \pt{Average value} \sh{7R2}: $f_\av$ = (total of $f$ over the region) divided by (the area of the region) = the height of the flat slab holding the same volume. \warn\ remember to divide. \ex $4x+2x^2y^2-9y^3$ on $[-4,0]\times[1,2]$: area 4; inner $[4xy+(2/3)x^2y^3-(9/4)y^4]\ev{1}{2}=4x+(14/3)x^2-135/4$; outer $[2x^2+(14/9)x^3-(135/4)x]\ev{-4}{0}=-(32-896/9+135)=-607/9$ \then\ $f_\av=-607/36$. \ex $4\cos x+7\cos y$ on $[0,\pi/2]\times[0,\pi]$: $\I{0}{\pi}\cos y\,dy=0$ and $\I{0}{\pi/2}4\cos x\,dx=4$ \then\ total $4\pi$, area $\pi^2/2$ \then\ $8/\pi$. \shc On a rectangle the average of $g(x)+h(y)$ = (1D average of $g$) + (1D average of $h$): $2\sin x+3\sin y$ on the same $R$: $4/\pi+6/\pi=10/\pi$. Midpoint estimate: $\Delta A$ × (sum of $f$ at the cell centers).

\hd{L24 DOUBLE INTEGRALS OVER GENERAL REGIONS}
\pt{Type I} \sh{7M2} (vertical strips, $dy\,dx$): the region lies between the vertical lines $x=a$ and $x=b$ (numbers); $g_1(x)$ = bottom curve, $g_2(x)$ = top curve = the inner limits. \pt{Type II} \sh{7M3} (horizontal strips, $dx\,dy$): between $y=c$ and $y=d$; $h_1(y)$ = left curve, $h_2(y)$ = right curve, both solved for $x$. $D$ = the region; $f$ = the integrand (height). \textbf{Limits recipe}: 1) sketch every boundary and shade the region; 2) intersections: set the curves equal; 3) slide a strip across: where it ENTERS the region is the lower inner limit, where it EXITS is the upper; 4) outer limits = first and last position of the strip. \warn\ outer limits are always numbers; inner limits may contain only the outer variable; a variable never shows up in a limit after it has been integrated away. If the entry or exit curve changes partway across, split the region or use the other type. \ex Prism: base = triangle bounded by the $x$-axis, $y=x$, $x=1$; top $z=3-x-y$. Type I: $\I{0}{1}\I{0}{x}(3-x-y)\,dy\,dx=\I{0}{1}(3x-(3/2)x^2)\,dx=3/2-1/2=1$. Type II: $\I{0}{1}\I{y}{1}(3-x-y)\,dx\,dy=\I{0}{1}(5/2-4y+(3/2)y^2)\,dy=1$ \ok. \ex $\II{D}(3x^2+y^2)\,dA$, $D$: $-2\le y\le3$, $y^2-3\le x\le y+3$ (already Type II: parabola on the left, line on the right): $\Ia{-2}{3}\Ia{y^2-3}{y+3}(3x^2+y^2)\,dx\,dy$; inner $[x^3+y^2x]$ \then\ $(y+3)^3+y^2(y+3)-(y^2-3)^3-y^2(y^2-3)=-y^6+8y^4+2y^3-12y^2+27y+54$; outer \then\ $2375/7$. \ex $\Ia{0}{1}\Ia{x^3}{x^3+7}(8x+5y)\,dy\,dx$: inner $56x+(5/2)((x^3+7)^2-x^6)=56x+35x^3+245/2$ \then\ $28+35/4+245/2=637/4$. \ex $\II{D}xy\,dA$, $D$ between $y=x^2$ and $y=x$: they meet where $x^2=x$, so $x=0,1$; in between the line is on top \then\ $\Ia{0}{1}\Ia{x^2}{x}xy\,dy\,dx=\I{0}{1}(x^3-x^5)/2\,dx=1/8-1/12=1/24$.
\pt{Switch the order} (when the inner integral cannot be done, like $e\cx(y^2)$ or $\sin(y^2)$, or is just uglier): sketch the region from the given limits, then describe the SAME region with the other kind of strip. Never just swap the differentials. \ex $\I{0}{1}\I{x}{1}e\cx(y^2)\,dy\,dx$: the region is $0\le x\le y\le1$ \then\ $\I{0}{1}\I{0}{y}e\cx(y^2)\,dx\,dy=\I{0}{1}ye\cx(y^2)\,dy=[(1/2)e\cx(y^2)]\ev{0}{1}=(e-1)/2$ ($u=y^2$). \ex $\I{0}{4}\I{-4}{y}f\,dx\,dy=\I{-4}{0}\I{0}{4}f\,dy\,dx+\I{0}{4}\I{x}{4}f\,dy\,dx$ (vertical strips change their bottom curve at $x=0$, so two pieces). \ex $\Ia{0}{4}\Ia{\rt x}{2}f\,dy\,dx=\Ia{0}{2}\Ia{0}{y^2}f\,dx\,dy$ (the curve $y=\rt x$, solved for $x$, is $x=y^2$).
\pt{Volume} = $\II{D}$(top $-$ bottom)$\,dA$ over the shadow $D$. ``First octant'' = $x,y,z\ge0$: solve the plane for $z$, set $z=0$ to get the triangle on the floor. \ex under $8x+2y+2z=32$: $z=16-4x-y$, floor edge $4x+y=16$ \then\ $\I{0}{4}\I{0}{16-4x}(16-4x-y)\,dy\,dx=\I{0}{4}(1/2)(16-4x)^2\,dx=512/3$ ($u=16-4x$). \shc A plane that cuts the axes at $a$, $b$, $c$ (set the other two variables to 0) cuts off first-octant volume $abc/6$: intercepts 4, 16, 16 \then\ $512/3$ \ok; $3x+y+2z=6$: $(2)(6)(3)/6=6$.
\pt{Area} \sh{7R1} = $\II{D}1\,dA$, which is Calc 1's top minus bottom. \ex between $y=1/(2x)$ and $y=2e^x$ from $x=1$ to $x=7$ ($2e^x$ is on top): $\I{1}{7}(2e^x-1/(2x))\,dx=[2e^x-(1/2)\ln x]\ev{1}{7}$ \sh{3R4}\,\sh{3R3} $=2e^7-2e-(1/2)\ln7$.
\pt{Average} \sh{7R2} = $(1/\text{Area of }D)\II{D}f\,dA$. \ex $D$ bounded by $y=0$, $x=y-7$, $x=7-y$ (triangle with corners $(-7,0)$, $(7,0)$, $(0,7)$; Type II), $f=9x+5y$: area 49; $\I{0}{7}\I{y-7}{7-y}(9x+5y)\,dx\,dy$: inner $[(9/2)x^2+5xy]=5y(14-2y)$ (the $9x$ part cancels by symmetry) \then\ $\I{0}{7}(70y-10y^2)\,dy=1715/3$ \then\ $f_\av=35/3$. \shc A linear $f=ax+by+c$ over a triangle averages to $f$ at the centroid = the mean of the three corner values: $(-63+63+35)/3=35/3$ \ok; $7x+4y$ on corners $(\pm14,0)$, $(0,14)$: $56/3$. \shc Terms odd in $x$ cancel over a region symmetric about the $y$-axis.

\hd{L25 DOUBLE INTEGRALS IN POLAR}
Use it for disks, rings, wedges, or whenever $x^2+y^2$ appears. $r$ = distance from the origin; $\theta$ = angle counterclockwise from the $+x$-axis, in radians. Replace $x=r\cos\theta$, $y=r\sin\theta$ \sh{5C2}, $x^2+y^2=r^2$ \sh{4L6} (the sheet's $f(r,\theta)$ = the integrand after this replacement) and \textbf{$dA=r\,dr\,d\theta$} (the extra $r$: a cell at radius $r$ is $dr$ long and $r\,d\theta$ wide).
\pt{Polar rectangle} \sh{7M4}: $a\le r\le b$ (inner and outer radius), $\alpha\le\theta\le\beta$ (first and last angle), all numbers. \pt{General polar region} \sh{7M5}: $h_1(\theta)$ = where a ray from the origin ENTERS the region, $h_2(\theta)$ = where it EXITS (0 to the curve when the origin is inside); $\alpha$, $\beta$ = first and last ray. \warn\ Count the $r$'s: a problem written with $dA$ needs the extra $r$ added; one already written with ``$r\,dr\,d\theta$'' has it. \textbf{$\theta$-ranges}: whole plane $0$ to $2\pi$; $y\ge0$: $0$ to $\pi$; $y\le0$: $\pi$ to $2\pi$; $x\ge0$: $-\pi/2$ to $\pi/2$; $x\le0$: $\pi/2$ to $3\pi/2$; quadrants I, II, III, IV: $0..\pi/2$, $\pi/2..\pi$, $\pi..3\pi/2$, $3\pi/2..2\pi$. Lines: $x=a$ is $r=a\sec\theta$; $y=b$ is $r=b\csc\theta$; $y=x$ is $\theta=\pi/4$; $y=mx$ is $\theta=\tan^{-1}m$ ($x>0$ half). \ex $\II{R}2\theta r\,dA$, $1\le r\le7$, $2\pi/3\le\theta\le3\pi/2$: $\Ia{2\pi/3}{3\pi/2}\Ia{1}{7}2\theta r^2\,dr\,d\theta$; inner $2\theta[r^3/3]\ev{1}{7}=228\theta$; outer $114[\theta^2]=114(9\pi^2/4-4\pi^2/9)=114\cdot65\pi^2/36=1235\pi^2/6$. \ex $\II{R}(x+y)\,dA$, $1\le x^2+y^2\le4$, $x\le0$: $\Ia{\pi/2}{3\pi/2}\Ia{1}{2}r^2(\cos\theta+\sin\theta)\,dr\,d\theta=(7/3)[\sin\theta-\cos\theta]$ \sh{3R6}\,\sh{3R8} $=(7/3)(-1-1)=-14/3$. \ex $\II{R}7y\,dA$, $36\le x^2+y^2\le81$, $y\le0$: $\I{\pi}{2\pi}\I{6}{9}7r^2\sin\theta\,dr\,d\theta=(7)(171)(-2)=-2394$. \shc Number limits and a product integrand: multiply the $r$-integral by the $\theta$-integral. \ex Cardioid: $\II{D}r^2\sin\theta\,r\,dr\,d\theta$, $D$ between the polar axis and the upper half of $r=1+\cos\theta$: $\I{0}{\pi}\I{0}{1+\cos\theta}r^3\sin\theta\,dr\,d\theta=\I{0}{\pi}(1/4)(1+\cos\theta)^4\sin\theta\,d\theta$; $u=1+\cos\theta$, $du=-\sin\theta\,d\theta$, limits $2\to0$ \then\ $(1/4)\I{0}{2}u^4\,du=8/5$.
\pt{Volume} = $\II{D}$(height)$\,r\,dr\,d\theta$; the base radius is where the surface meets the floor (set $z=0$) or where two surfaces meet. \ex solid bounded by $z=2-9x^2-9y^2$ and the planes $x=0$, $y=0$, $z=0$ (the first-quadrant quarter): $2-9r^2=0$ \then\ $r=\rt2/3$: $\Ia{0}{\pi/2}\Ia{0}{\rt2/3}(2-9r^2)\,r\,dr\,d\theta=(\pi/2)[r^2-(9/4)r^4]=(\pi/2)(2/9-1/9)=\pi/18$. \ex under the cone $z=\rt(x^2+y^2)$, above the disk $3x^2+3y^2\le4$: $z=r$, $r\le2/\rt3$: $(2\pi)(1/3)(2/\rt3)^3=16\rt3\,\pi/27$.
\pt{Area} \sh{7R3}: $\alpha,\beta$ = angle range; $h_1(\theta),h_2(\theta)$ = entry and exit radius. \shc The $r$-integral is always $(1/2)$(exit$^2-$ entry$^2$), so $A=\J(1/2)$(exit$^2-$ entry$^2$)$\,d\theta$ (entry 0: this is \sh{4L7}). Squaring the curve gives $\sin^2$ or $\cos^2$ \then\ power-reducing \sh{p1}: $\sin^2u=(1-\cos2u)/2$, $\cos^2u=(1+\cos2u)/2$. \ex $r=-6\sin4\theta$, $0\le\theta\le3\pi/2$: $\I{0}{3\pi/2}18\sin^2(4\theta)\,d\theta=\I{0}{3\pi/2}(9-9\cos8\theta)\,d\theta=[9\theta-(9/8)\sin8\theta]=27\pi/2$ ($\sin12\pi=0$). A negative $r$ only draws the point on the opposite ray; the square makes its area count positive. \ex quadrant-I petal of $r=8\sin2\theta$: $\I{0}{\pi/2}32\sin^2(2\theta)\,d\theta=(32)(\pi/4)=8\pi$. \ex ring wedge $3\le r\le7$, $0\le\theta\le\pi/2$: $(1/2)(49-9)(\pi/2)=10\pi$. \shc Roses $r=a\cos n\theta$ or $a\sin n\theta$: $n$ odd \then\ $n$ petals (all traced by $0..\pi$); $n$ even \then\ $2n$ petals ($0..2\pi$); \textbf{each petal has area $\pi a^2/(4n)$}, so count the petals your $\theta$-range traces (above: 6 of the 8 petals × $9\pi/4$ \ok). Other curves: $r=a$ circle; $r=2a\cos\theta$ circle through the origin centered $(a,0)$, $-\pi/2..\pi/2$; $r=2a\sin\theta$ centered $(0,a)$, $0..\pi$; cardioid $r=a(1\pm\cos\theta)$, area $3\pi a^2/2$. Curves cross where they are equal: $14\cos\theta=7$ \then\ $\theta=\pm\pi/3$. \ex disk inside $r=2\cos\theta$: $\Ia{-\pi/2}{\pi/2}\Ia{0}{2\cos\theta}r\,dr\,d\theta=\Ia{-\pi/2}{\pi/2}2\cos^2\theta\,d\theta=\pi$. \ex $r=\sin\theta-2$ ($r<0$ throughout): whole loop $(1/2)\I{0}{2\pi}(\sin\theta-2)^2\,d\theta=(1/2)(\pi+8\pi)=9\pi/2$.
\pt{Convert a Cartesian integral}: read the region off the limits, sketch it, describe it again with $r$ and $\theta$, convert the integrand, add the $r$. \ex $\Ia{0}{2}\Ia{0}{\rt(4-x^2)}(x^2+y^2)\,dy\,dx$ = quarter disk of radius 2 \then\ $\I{0}{\pi/2}\I{0}{2}r^3\,dr\,d\theta=2\pi$. \ex $\II{}e\cx(-x^2-y^2)\,dA$ over $x^2+y^2\le1$: $2\pi\I{0}{1}re\cx(-r^2)\,dr=\pi(1-1/e)$ ($u=r^2$). \warn\ $\rt(x^2+y^2)=r$, not $r^2$.

\hd{L26 TRIPLE INTEGRALS}
\pt{Meaning.} $\III{E}f\,dV$: chop the solid $E$ into tiny boxes of volume $dV=dx\,dy\,dz$ and add up (value of $f$ there) × (tiny volume). $f=1$ \then\ volume \sh{7R5}; $f$ = density \then\ mass; $f$ = charge density \then\ total charge; divide by the volume \then\ average \sh{7R4}. \pt{Box} \sh{7M6}: $B=[a,b]\times[c,d]\times[e,f]$ = the $x$, $y$, $z$ ranges (the limit ``$f$'' there is a number, not the function); any of the 6 orders gives the same value; do the innermost differential first with the other letters frozen. \ex $\I{0}{1}\I{2}{4}\I{1}{5}(x+yz^2)\,dx\,dy\,dz$: $x$: $[x^2/2+xyz^2]\ev{1}{5}=12+4yz^2$; $y$: $[12y+2y^2z^2]\ev{2}{4}=24+24z^2$; $z$: $[24z+8z^3]\ev{0}{1}=32$. \shc Number limits and integrand $g(x)h(y)k(z)$ \then\ multiply three single integrals. \ex $\I{-3}{-2}\I{1}{2}\I{0}{\ln3}y^2z^3e\cx(2x)\,dx\,dy\,dz=[e\cx(2x)/2]\cdot[y^3/3]\cdot[z^4/4]$ \sh{3R5}\,\sh{3R2} $=(4)(7/3)(-65/4)=-455/3$ (using $e\cx(2\ln3)=9$). \ex $\I{1}{2}\I{2}{3}\I{1}{3}7yz^2/(5x^3)\,dx\,dy\,dz=(7/5)(4/9)(5/2)(7/3)=98/27$.
\pt{General solid} \sh{7M7} (order $dz\,dy\,dx$): $u_1(x,y)$, $u_2(x,y)$ = the bottom and top surfaces, solved for $z$; $g_1(x)$, $g_2(x)$ = the two edges of the shadow $D$ on the $xy$-plane; $a,b$ = the $x$-range. \how 1) push a vertical $z$-stick through the solid: it enters at $u_1$ and leaves at $u_2$; 2) flatten the solid to its shadow (set $z=0$ in the top surface, or intersect top with bottom); a $y$-line across the shadow enters at $g_1(x)$ and leaves at $g_2(x)$; 3) $x$ runs from $a$ to $b$. A limit may contain only variables that are still waiting to be integrated. \shc No $z$ in the integrand \then\ the $z$-integral just multiplies by the stick length (top $-$ bottom). \ex $\III{E}(5x-3y)\,dV$ over the tetrahedron bounded by $x=0$, $y=0$, $z=0$, $x+y+z=1$: $\I{0}{1}\I{0}{1-x}\I{0}{1-x-y}(5x-3y)\,dz\,dy\,dx$; $z$: $(5x-3y)(1-x-y)$; $y$, with the nickname $a=1-x$ (a constant during this step): $\I{0}{a}(5x-3y)(a-y)\,dy=(5/2)xa^2-a^3/2$; $x$, with $u=1-x$ (so $x=1-u$, $dx=-du$, limits $1\to0$ flip back to $0\to1$): $\I{0}{1}[(5/2)(1-u)u^2-u^3/2]\,du=5/6-5/8-1/8=1/12$. \ex set-up only: $9x+4y$ over the tetrahedron cut off by $4x+7y+9z=252$: top $z=28-4x/9-7y/9$; $z=0$ \then\ $y=36-4x/7$; then $y=0$ \then\ $x=63$: $\I{0}{63}\I{0}{36-4x/7}\I{0}{28-4x/9-7y/9}(9x+4y)\,dz\,dy\,dx$.
\pt{Volume} \sh{7R5} = $\III{E}1\,dV$. \ex below $z=4x^2+4y^2+4$, above the $xy$-plane, shadow $[0,3]\times[0,2]$: $\I{0}{3}\I{0}{2}(4x^2+4y^2+4)\,dy\,dx=\I{0}{3}(8x^2+56/3)\,dx=72+56=128$. \shc First-octant tetrahedron under a plane with intercepts $a,b,c$: $V=abc/6$; for $Ax+By+Cz=D$ that is $D^3/(6ABC)$: $8x+y+4z=32$ \then\ $(4)(32)(8)/6=512/3$. \warn\ ``$y\ge0$, $z\ge0$'' leaves $x$ free: under $z=16-x^2-3y^2$ with $y,z\ge0$ the limits are $-4\le x\le4$, $0\le y\le\rt((16-x^2)/3)$, $0\le z\le16-x^2-3y^2$ ($V=64\pi/\rt3$).
\pt{Average} \sh{7R4} = $(1/V(E))\III{E}f\,dV$, where $V(E)$ = volume of the solid. \ex temperature $T=xy+8x+20$ over the tetrahedron $x+y+z\le1$ in the first octant: $V=1/6$; $z$: $(xy+8x+20)(1-x-y)$; $y$ ($a=1-x$): $xa^3/6+(8x+20)a^2/2$; $x$ ($u=1-x$): $\I{0}{1}[(1-u)u^3/6+(14-4u)u^2]\,du=1/120+11/3=147/40$ \then\ $T_\av=(6)(147/40)=441/20=22.05$. \shc Over that corner tetrahedron $\III{}x^py^qz^s\,dV=p!\,q!\,s!/(p+q+s+3)!$ (triangle $x+y\le1$: $\II{}x^py^q\,dA=p!\,q!/(p+q+2)!$): $\III{}xy=1/120$, $\III{}x=1/24$, $\III{}1=1/6$ \then\ $1/120+8/24+20/6=147/40$ \ok. \shc On a box the average of a product $g(x)h(y)k(z)$ = the product of the three 1D averages: $x^2yz^2$ on the cube $[0,6]^3$: $(12)(3)(12)=432$ (the average of $x^n$ on $[0,L]$ is $L^n/(n+1)$). \shc Integrand odd in $z$, solid symmetric in $z$: 0.

\hd{L27 TRIPLE INTEGRALS: CYLINDRICAL \& SPHERICAL}
\pt{Which coordinates?} Flat faces, boxes, tetrahedra \then\ rectangular. $x^2+y^2$, cylinders, paraboloids, a round shadow \then\ cylindrical. $x^2+y^2+z^2$, balls, cones with tip at the origin \then\ spherical. Three jobs every time: convert the integrand, turn the region into limits, and include the volume factor ($r$ or $\rho^2\sin\phi$) exactly once.
\pt{Cylindrical} $(r,\theta,z)$ = polar coordinates on the floor plus the ordinary height: \sh{5C1}\,\sh{5C2}; $r$ = distance from the $z$-axis; $dV=r\,dz\,dr\,d\theta$. Surfaces: $r=c$ cylinder around the $z$-axis; $\theta=c$ vertical half-plane hinged on the $z$-axis; $z=c$ horizontal plane; cone $z^2=c^2(x^2+y^2)$ is $z=\pm cr$; sphere $x^2+y^2+z^2=c^2$ is $r^2+z^2=c^2$; paraboloid $z=c(x^2+y^2)$ is $z=cr^2$. \textbf{Box} \sh{7M8}: $a\le r\le b$, $\alpha\le\theta\le\beta$, $c\le z\le d$ (the sheet's $g(x,y,z)$ is the original integrand and $f(r,\theta,z)$ is the same thing after substituting). \textbf{General} \sh{7M9}: $u_1(r,\theta)$, $u_2(r,\theta)$ = bottom and top surfaces written with $r,\theta$ (the $z$-stick); $g_1(\theta)$, $g_2(\theta)$ = where a ray enters and leaves the shadow; $\alpha,\beta$ = angle range. \ex $\III{B}(zr\sin\theta)\,r\,dr\,d\theta\,dz$, $0\le r\le2$, $0\le\theta\le\pi/2$, $0\le z\le4$ (the second $r$ is the volume factor, already included): $(\I{0}{2}r^2\,dr)(\I{0}{\pi/2}\sin\theta\,d\theta)(\I{0}{4}z\,dz)=(8/3)(1)(8)=64/3$. \ex $\III{E}1/(x+2)\,dV$, $E$: $x^2+y^2\le4$, $x\ge0$, $y\ge0$, $0\le z\le x+2$: $\I{0}{\pi/2}\I{0}{2}\I{0}{r\cos\theta+2}r/(r\cos\theta+2)\,dz\,dr\,d\theta$; the stick length cancels the denominator \then\ $\I{0}{\pi/2}\I{0}{2}r\,dr\,d\theta=\pi$. \ex $5x^2+5y^2$ over the same quarter disk with $0\le z\le2-x$: $\I{0}{2}\I{0}{\pi/2}5r^3(2-r\cos\theta)\,d\theta\,dr=\I{0}{2}5r^3(\pi-r)\,dr=20\pi-32$. \ex between the paraboloid $z=x^2+y^2$ and the plane $z=4$: they meet where $r^2=4$ \then\ $\Ia{0}{2\pi}\Ia{0}{2}\Ia{r^2}{4}r\,dz\,dr\,d\theta=2\pi\I{0}{2}(4r-r^3)\,dr=8\pi$. \ex above the cone $z=r$, inside the sphere $r^2+z^2=8$: they meet where $2r^2=8$ \then\ $\Ia{0}{2\pi}\Ia{0}{2}\Ia{r}{\rt(8-r^2)}r\,dz\,dr\,d\theta=2\pi\I{0}{2}(r\rt(8-r^2)-r^2)\,dr=(32\pi/3)(\rt2-1)$ ($u=8-r^2$).
\pt{Spherical} $(\rho,\theta,\phi)$ \sh{5C3}\,\sh{5C4}: $\rho$ = distance from the origin; $\theta$ = the same around-the-$z$-axis angle as in polar; $\phi$ = angle measured DOWN from the $+z$-axis, $0\le\phi\le\pi$ (0 = straight up, $\pi/2$ = the $xy$-plane, $\pi$ = straight down). $z=\rho\cos\phi$, $r=\rho\sin\phi$ \sh{5C5}, $x^2+y^2+z^2=\rho^2$. $dV=\rho^2\sin\phi\,d\rho\,d\phi\,d\theta$ (a brick with edges $d\rho$, $\rho\,d\phi$, $\rho\sin\phi\,d\theta$). \textbf{Box} \sh{7M10}: $a\le\rho\le b$, $\alpha\le\theta\le\beta$, $\gamma\le\phi\le\psi$. \textbf{General} \sh{7M11} (\warn\ printed in the order $d\phi\,d\rho\,d\theta$, with $\phi$ innermost between $u_1(\rho,\theta)$ and $u_2(\rho,\theta)$ and $\rho$ between $g_1(\theta)$ and $g_2(\theta)$; the usual order is $d\rho\,d\phi\,d\theta$, and any order is fine as long as each limit uses only variables further out). \textbf{Region into limits}: inside the sphere of radius $R$: $0\le\rho\le R$; $z\ge0$: $0\le\phi\le\pi/2$; $z\le0$: $\pi/2\le\phi\le\pi$; $y\ge0$: $0\le\theta\le\pi$; $x\ge0$: $-\pi/2\le\theta\le\pi/2$; first octant: both angles $0$ to $\pi/2$; whole ball: $\theta$ $0$ to $2\pi$, $\phi$ $0$ to $\pi$. \warn\ $z$ decides $\phi$; $x$ and $y$ decide $\theta$. Cone $z=\rt(x^2+y^2)$ is $\phi=\pi/4$ (in general $\tan\phi=r/z$); plane $z=c$ is $\rho=c\sec\phi$; sphere $x^2+y^2+z^2=2az$ is $\rho=2a\cos\phi$. \ex $\I{0}{2\pi}\I{0}{\pi/2}\I{0}{1}\rho^2\sin\phi\,d\rho\,d\phi\,d\theta=(1/3)(1)(2\pi)=2\pi/3$ (half the unit ball \ok). \ex $5x^2+5y^2+5z^2$ inside $x^2+y^2+z^2=1$ with $z\ge0$, $y\ge0$: $f=5\rho^2$ \then\ $\I{0}{\pi}\I{0}{\pi/2}\I{0}{1}5\rho^4\sin\phi\,d\rho\,d\phi\,d\theta=(1)(1)(\pi)=\pi$. \ex set up $f=-6$ over $x^2+y^2+z^2\le33$, $z\ge0$: $\Ia{0}{2\pi}\Ia{0}{\pi/2}\Ia{0}{\rt33}(-6)\rho^2\sin\phi\,d\rho\,d\phi\,d\theta$ (value $-132\rt33\,\pi$). \ex $f=z$ over the upper half-ball of radius 2: $z=\rho\cos\phi$ \then\ $\I{0}{2\pi}\I{0}{\pi/2}\I{0}{2}\rho^3\cos\phi\sin\phi\,d\rho\,d\phi\,d\theta=(2\pi)(1/2)(4)=4\pi$. \shc Constant $f$ \then\ $f$ × volume: $f=9$ on the lower half-ball of radius 3: $(9)(2/3)(\pi)(27)=162\pi$. \ex Factor before converting: $4x^2z+4y^2z+4z^3=4z(x^2+y^2+z^2)=4\rho^3\cos\phi$ over $x^2+y^2+z^2\le169/4$ with $y,z\ge0$ ($x$ is free, so $0\le\theta\le\pi$): $\I{0}{\pi}\I{0}{\pi/2}\I{0}{13/2}4\rho^5\cos\phi\sin\phi\,d\rho\,d\phi\,d\theta=(\pi)(1/2)(4/6)(13/2)^6=4826809\pi/192$. \shc $\I{0}{\pi/2}\sin\phi\,d\phi=1$; $\I{0}{\pi}\sin\phi\,d\phi=2$; $\I{0}{\pi/2}\sin\phi\cos\phi\,d\phi=1/2$. \shc Ice-cream cone (ball of radius $R$ inside the cone $\phi\le\phi_0$): $V=(2\pi/3)R^3(1-\cos\phi_0)$; inside $\rho\le2$ above the cone $z=\rt(x^2+y^2)$: $(2\pi/3)(8)(1-\rt2/2)=(8\pi/3)(2-\rt2)$. Sanity volumes: ball $(4/3)\pi R^3$, cylinder $\pi R^2h$, cone $(1/3)\pi R^2h$, shell $a\le\rho\le b$: $(4/3)\pi(b^3-a^3)$.

\hd{BY-HAND TOOLKIT · RULES THE SHEET DOES NOT GIVE}
\pt{Partial derivatives} (used in every section): $f_x=\pd f/\pd x$ = differentiate in $x$ while every other letter is treated as a constant number. A frozen factor stays, $\pd_x(3xy)=3y$; a frozen term dies, $\pd_x(2y^2)=0$. Chain rule: $\pd_xe\cx(xy)=ye\cx(xy)$, $\pd_y\sin(x^2y)=x^2\cos(x^2y)$. Product rule only when both factors contain the variable. Second partials: $f_{xx}=(f_x)_x$, $f_{yy}=(f_y)_y$, $f_{xy}=(f_x)_y=f_{yx}$.
\pt{Derivative rules.} On the sheet, \sh{p3} left column, rows 1–16: $k$, $x^n$, $\ln\abs{x}$, $e^x$, $e^{bx}$, $\sin x$, $\sin bx$, $\cos x$, $\cos bx$, $\tan$, $\csc$, $\sec$, $\cot$, $\sin^{-1}$, $\tan^{-1}$, $\sec^{-1}$. NOT on the sheet: \textbf{product} $(fg)'=f'g+fg'$; \textbf{quotient} $(f/g)'=(f'g-fg')/g^2$ (bottom times derivative of top, minus top times derivative of bottom, over bottom squared); \textbf{chain} $[f(g(x))]'=f'(g(x))\cdot g'(x)$ = derivative of the outside with the inside left untouched, times the derivative of the inside. Chain forms ($u$ = the inside, $u'$ = its derivative): $(u^n)'=nu^{n-1}u'$; $(\rt u)'=u'/(2\rt u)$; $(1/u)'=-u'/u^2$; $(e^u)'=u'e^u$; $(a^u)'=u'a^u\ln a$; $(\ln u)'=u'/u$; $(\sin u)'=u'\cos u$; $(\cos u)'=-u'\sin u$; $(\tan u)'=u'\sec^2u$; $(\tan^{-1}u)'=u'/(1+u^2)$. Rewrite before differentiating: $\rt x=x^{1/2}$, $1/x^3=x^{-3}$, cube root of $t$ = $t^{1/3}$. Exponents: $t^at^b=t^{a+b}$ (so $t^{1/3}\cdot t=t^{4/3}$), $t^a/t^b=t^{a-b}$ ($t^2/t^{3/4}=t^{5/4}$), $(t^a)^b=t^{ab}$.
\pt{Integral rules.} On the sheet, \sh{p3} right column, same rows: $k$; $x^n$ ($n\ne-1$); $1/x\to\ln\abs{x}$; $e^x$; $e^{bx}\to e^{bx}/b$; $\cos x$; $\cos bx\to(\sin bx)/b$; $\sin x$; $\sin bx\to-(\cos bx)/b$; $\sec^2$; $\csc\cot$; $\sec\tan$; $\csc^2$; $1/\rt(1-x^2)$; $1/(1+x^2)$; $1/(x\rt(x^2-1))$. NOT on the sheet: $\J\tan u=-\ln\abs{\cos u}$; $\J\cot u=\ln\abs{\sin u}$; $\J\sec u=\ln\abs{\sec u+\tan u}$; $\J\ln x=x\ln x-x$; $\J a^x=a^x/\ln a$; $\J(ax+b)^n=(ax+b)^{n+1}/(a(n+1))$; $\J1/(ax+b)=(1/a)\ln\abs{ax+b}$; $\J1/(a^2+x^2)=(1/a)\tan^{-1}(x/a)$; $\J\sin^2u=u/2-(\sin2u)/4$; $\J\cos^2u=u/2+(\sin2u)/4$ (both from power-reducing \sh{p1}); $\J\sin u\cos u=(1/2)\sin^2u$. Definite: $\I{a}{b}f=F(b)-F(a)$, top limit minus bottom limit (\warn\ $F$ at 0 is not always 0: $\cos0=1$, $e^0=1$). Constants come out front; integrate term by term. Handy values: $\sin^2$ or $\cos^2$ integrated over $0..\pi/2$ gives $\pi/4$, over $0..\pi$ gives $\pi/2$, over $0..2\pi$ gives $\pi$; $\sin$ and $\cos$ over a full period give 0. Symmetric limits: an odd integrand gives $\I{-a}{a}=0$; an even one gives $2\I{0}{a}$.
\pt{u-substitution} = the chain rule run backwards. Use it when the integrand is (a function of some inside chunk) × (the derivative of that chunk, give or take a constant factor). \how 1) $u$ = the inside chunk (what sits inside a power, a root, an exponent, a trig function or a denominator); 2) differentiate: $du=u'\,dx$, so $dx=du/u'$; 3) substitute: every $x$ has to disappear (the leftover factor cancels against $u'$; a missing CONSTANT is fine: put $1/c$ out front); 4) integrate in $u$; 5) indefinite: put the $x$-expression back. Definite: convert the limits ($x=a$ becomes $u(a)$, $x=b$ becomes $u(b)$) and never go back to $x$, OR put $x$ back and use the old limits. Never mix the two. \ex $\J y^2\cos(6y^3)\,dy$: $u=6y^3$, $du=18y^2\,dy$, so $y^2\,dy=du/18$ \then\ $(1/18)\J\cos u\,du=(1/18)\sin(6y^3)$. \ex $\I{0}{4}(1/2)(16-4x)^2\,dx$: $u=16-4x$, $dx=-du/4$, limits $16\to0$ \then\ $-(1/8)\I{16}{0}u^2\,du=(1/8)\I{0}{16}u^2\,du=16^3/24=512/3$ (a minus sign flips the limits). \ex $\I{0}{\pi}\sin\theta\,(1+\cos\theta)^4\,d\theta$: $u=1+\cos\theta$, $du=-\sin\theta\,d\theta$, limits $2\to0$ \then\ $\I{0}{2}u^4\,du=32/5$. \ex Leftover $x$: $\I{0}{1}x(1-x)^3\,dx$: $u=1-x$, so $x=1-u$, $dx=-du$, limits $1\to0$ \then\ $\I{0}{1}(1-u)u^3\,du=1/4-1/5=1/20$. \ex $\I{0}{2}r\rt(4-r^2)\,dr$: $u=4-r^2$, $r\,dr=-du/2$, limits $4\to0$ \then\ $(1/2)\I{0}{4}\rt u\,du=[(1/3)u^{3/2}]\ev{0}{4}=8/3$. \ex $\I{0}{1}re\cx(-r^2)\,dr$: $u=-r^2$, $r\,dr=-du/2$ \then\ $[-(1/2)e\cx(-r^2)]\ev{0}{1}=(1-1/e)/2$. \ex $\J x/(x^2+1)\,dx$: $u=x^2+1$ \then\ $(1/2)\ln(x^2+1)$. \ex Trig inside: $\J\cos^2\phi\sin\phi\,d\phi$: $u=\cos\phi$, $du=-\sin\phi\,d\phi$ \then\ $-(1/3)\cos^3\phi$. \ex Odd power: $\J\sin^3\phi\,d\phi=\J(1-\cos^2\phi)\sin\phi\,d\phi=-\cos\phi+(1/3)\cos^3\phi$ (from 0 to $\pi$: $4/3$). \shc \textbf{Linear inside}: integrate as if the inside were plain $x$, then divide by its coefficient: $\J\cos(6y)=(\sin6y)/6$; $\J e\cx(x-7)=e\cx(x-7)$; $\J\rt(6x+7)=(6x+7)^{3/2}/9$; $\J(16-4x)^2=-(16-4x)^3/12$. \shc Guess and check: write down the outside's antiderivative, differentiate it, fix the constant. \warn\ If the missing factor contains the variable (as in $\J e\cx(x^2)\,dx$, which has no spare $x$) u-sub fails: in a double integral switch the order (L24); a product like $xe^x$ needs parts. A nickname such as $a=1-x$ used inside a $y$-integral is NOT a substitution (it is a constant there): no $du$, no new limits.
\pt{Integration by parts} (a product of unlike types): $\J u\,dv=uv-\J v\,du$; choose $u$ by LIATE (Logs, Inverse trig, Algebraic, Trig, Exponential), $dv$ = the rest. $\J xe^x=xe^x-e^x$; $\J x\cos x=x\sin x+\cos x$; $\J x\sin x=-x\cos x+\sin x$; $\J x\ln x=(x^2/2)\ln x-x^2/4$.
\pt{Trig.} Unit circle \sh{p2}: the point at angle $\theta$ is $(\cos\theta,\sin\theta)$. $\sin k\pi=0$, $\cos k\pi=(-1)^k$; $\sin(\pi/2)=1$, $\sin(3\pi/2)=-1$, $\cos(\pi/2)=\cos(3\pi/2)=0$; subtract $2\pi$ from large angles ($9\pi/4\to\pi/4$). $\pi/6$: $(\rt3/2,1/2)$; $\pi/4$: $(\rt2/2,\rt2/2)$; $\pi/3$: $(1/2,\rt3/2)$. Identities \sh{p1}: $\sin^2+\cos^2=1$; $\sin2\theta=2\sin\theta\cos\theta$; power-reducing (see L25). $\sec=1/\cos$, $\csc=1/\sin$.
\pt{Algebra.} Two linear equations: add or subtract a multiple of one to kill a variable, then back-substitute; or Cramer: $ax+by=p$, $cx+dy=q$ \then\ $x=(pd-bq)/(ad-bc)$, $y=(aq-pc)/(ad-bc)$. Quadratic formula $x=(-b\pm\rt(b^2-4ac))/(2a)$; parabola vertex $x=-b/(2a)$. Complete the square: $x^2+bx=(x+b/2)^2-b^2/4$. $(a\pm b)^2=a^2\pm2ab+b^2$; $(a+b)^3=a^3+3a^2b+3ab^2+b^3$. Roots: $\rt(a^2b)=a\rt b$ ($\rt1160=2\rt290$); rationalize $a/\rt b=a\rt b/b$. Vectors: length, unit vector, dot product \sh{4L11}\,\sh{4L12}\,\sh{4L16}. Logs: $\ln1=0$, $\ln e=1$, $e\cx(\ln a)=a$, $e\cx(2\ln3)=9$, $\ln a-\ln b=\ln(a/b)$.
\pt{Answer checks.} Direction vector really length 1? Gradient evaluated before the dot product? Implicit: minus sign, right letter on the bottom? $D$ checked at every critical point? Every Lagrange branch, plus the interior for a filled region? Inner limits use only outer variables, lower below upper? Extra $r$ or $\rho^2\sin\phi$ included exactly once? Average divided by the area or volume? Area and volume positive? Exact form ($\pi$, radicals) and radians? Did they ask for a point, a value or a vector? Inner integral done first? Trig values at the limits right?

\vfill
% hollow letters: white fill + black stroke (change [1] to e.g. [0.85] for a light gray fill)
\noindent\hfil\resizebox{!}{\dimexpr\stamplines\pitch-3.4pt\relax}{\vbox{\special{pdf:bc [1] [0]}\stampf\fontsize{20}{20}\selectfont
  \baselineskip=0pt \lineskiplimit=0pt \lineskip=2.6pt
  \halign{\hfil#\hfil\cr STAMP\cr HERE\cr}\special{pdf:ec}}}\hfil\par\kern1pt

\end{multicols*}
\end{document}
